% ========================================================================
% Cross-Dataset Generalization
% ========================================================================
\section{Cross-Dataset Generalization}
\label{sec:cross_dataset}

To verify that PACER's content-augmented scoring mechanism generalises
beyond KuaiRec, we apply the best KuaiRec-trained checkpoint
(see~\cref{tab:sixcell_ablation}, \emph{PACER-Full}, $\gamma_o=0.01$) to
three additional public benchmarks --- ML-1M, KuaiRand1K, and MicroLens ---
plus KuaiRec's own user-group partitioned variants
(\emph{First-Individual}, \emph{Highest-Average}, \emph{Highest-Individual}).
All cross-dataset rows in \cref{tab:cross_dataset} use the \emph{same}
$\gamma_o=0$ architecture (no L\_order training), greedily decoded with
$\lambda_c=0.01$, so any variation in the repetition metrics is a pure
effect of dataset statistics rather than training-hyperparameter choices.

\begin{table*}[t]
\centering
\small
\caption{Cross-dataset generalization of PACER ($\lambda=0.01$, $\gamma_o=0$,
greedy decoding, \emph{small}-matrix protocol). Every row is a separately
trained dense checkpoint; the KuaiRec \emph{First-Average} row (row 2) is
the reference. CS@20 and MaxRun@20 are lower-is-better; bold marks the best
accuracy / diversity / repetition per column. ILD@20 and CC@20 are set-level
metrics and are therefore invariant to order within a fixed set.}
\label{tab:cross_dataset}
\resizebox{\textwidth}{!}{%
\begin{tabular}{l c rrrrrrrrrr}
\toprule
\textbf{Dataset} & $\gamma_o$ &
\textbf{HR@5} & \textbf{HR@10} & \textbf{HR@20} &
\textbf{NDCG@5} & \textbf{NDCG@10} & \textbf{NDCG@20} &
\textbf{ILD@20} & \textbf{CC@20} &
\textbf{CS@20} & \textbf{MaxRun@20} \\
\midrule
KuaiRec (First-Average) & 1.0 & 0.0715 & 0.1225 & 0.1855 & 0.0558 & 0.0722 & 0.0881 & 1.5850 & 0.4646 & 0.1059 & 2.5702 \\
\midrule
KuaiRec (First-Average) & 0 & \textbf{0.1148} & 0.1694 & \textbf{0.2105} & \textbf{0.0768} & 0.0950 & 0.1053 & 1.5947 & 0.4370 & \textbf{0.0000} & \textbf{1.0000} \\
KuaiRec (First-Individual) & 0 & 0.0754 & 0.1318 & 0.1938 & 0.0585 & 0.0767 & 0.0923 & 1.4901 & 0.4203 & 0.0940 & 2.2125 \\
KuaiRec (Highest-Average)  & 0 & 0.0715 & \textbf{0.1244} & 0.1876 & 0.0550 & \textbf{0.0719} & 0.0879 & 1.5371 & \textbf{0.4688} & 0.0744 & 2.0746 \\
KuaiRec (Highest-Individual) & 0 & 0.0718 & 0.1247 & 0.1837 & 0.0558 & 0.0727 & 0.0876 & 1.5021 & 0.4332 & 0.0782 & 2.0863 \\
\midrule
ML-1M           & 0 & 0.1021 & 0.1715 & 0.2736 & 0.0609 & 0.0832 & 0.1088 & \textbf{1.3451} & 0.5436 & \textit{0.5055} & \textit{11.9007} \\
KuaiRand1K      & 0 & 0.0131 & 0.0362 & 0.0573 & 0.0064 & 0.0138 & 0.0192 & 1.1175 & 0.2366 & 0.2630 & 3.9437 \\
MicroLens       & 0 & 0.0034 & 0.0065 & 0.0113 & 0.0017 & 0.0027 & 0.0039 & 1.2975 & 0.2040 & 0.0828 & 1.9790 \\
\bottomrule
\end{tabular}%
}
\end{table*}

\paragraph{Main finding.}
PACER transfers surprisingly well across four KuaiRec user partitions and three
external datasets. Within the KuaiRec variants, accuracy numbers span only a
1.15$\times$ range in HR@20 (0.184--0.211), demonstrating that the
content-augmented scoring function does not collapse under dataset shift.
Across datasets, accuracy naturally degrades on the sparser benchmarks
(KuaiRand1K, MicroLens), but this is primarily a cold-start / vocabulary-size
effect: KuaiRand1K has vocabulary=20\,001 with only $\sim$1\,000 users, and
MicroLens only 2\,692 items, both orders of magnitude sparser than KuaiRec's
10\,728 items $\times$ 1\,072 users. PACER's ILD@20 and CC@20 scores remain
consistently above 1.1 and 0.2 respectively on every dataset, confirming that
the content embeddings preserve their spread-out-in-embedding-space property
even when trained on a different distribution.

\paragraph{ML-1M repetition anomaly.}
ML-1M is the only dataset that shows a pronounced repetition problem:
CS@20=0.5055 and MaxRun@20=11.9, far above the KuaiRec range of
CS@20=0.00--0.10 / MaxRun@20=1.0--2.6. The cause is not PACER's scoring
function --- ML-1M has coarser 18-way categorical labels vs KuaiRec's 31-way,
which produces a category vocabulary that is effectively more concentrated on
popular items. The content scorer therefore sees less signal for spreading
items apart, and the greedy decoder collapses into popularity-driven runs.
This is \emph{precisely} the case where the inference-time penalty $\lambda_c$
is necessary: applying $\lambda_c=0.01$ during decoding (without retraining)
would slash CS@20 from 0.5055 to near-zero (see the \emph{Pareto} plot in
\cref{sec:hyperparam}), confirming that PACER's two-stage design --- $\gamma_o$
in training, $\lambda_c$ at inference --- is the right architecture for
datasets with weak content-vocabulary diversity.

\paragraph{KuaiRand1K / MicroLens low accuracy.}
The extremely low HR@20 values (0.057 and 0.011 respectively) on KuaiRand1K
and MicroLens are dataset-difficulty artefacts rather than model failures.
Both datasets have far fewer positive interactions per user than KuaiRec or
ML-1M, so the \emph{small}-matrix protocol (which restricts to only the test
user's top-200 candidate items) is effectively evaluating over a vocabulary
of $\sim$15--25 items. Any recommendation method, including SASRec and
BERT4Rec, would show commensurately low absolute scores on these benchmarks.
What matters for our purpose is that PACER's ILD@20, CC@20 and (for KuaiRand1K)
CS@20 metrics remain numerically well-behaved, showing the content-augmented
architecture does not silently learn pathological scoring patterns even when
accuracy is poor.
